\subsubsection{CenterCrop} \label{CenterCrop}
\begin{figure}
	 \centerline{\includegraphics[width=1\textwidth]{src/augmentations/CenterCrop.png}}
	 {\caption{Example of CenterCrop with proportion of 0.2 and 0.4 }\label{fig:CenterCrop}}
\end{figure}

CenterCrop transformation is a type of image augmentation technique that involves cropping an image from its center to generate a smaller sub-image. This transformation can be used to remove irrelevant areas from the image and focus on the most important features. The CenterCrop transformation can also help to standardize the size of images used for training machine learning models. The CenterCrop transformation can be advantageous for aerial images, as it can eliminate irrelevant background areas, enabling the focus on essential features. This is particularly useful for our dataset, which contains numerous images with center biases. 

However, it is crucial to note that in some pictures where there is no center bias, such as the Kormoran image in Figure \ref{fig:CenterCrop} or images showing a flock of birds, this transformation may remove important information from the image. Therefore, it is necessary to evaluate each image carefully and determine whether the CenterCrop transformation is suitable, as its usage may depend on the image's characteristics.

In our experiment, we will apply CenterCrop transformation to our aerial images, using proportions of 0.2 and 0.4, as shown in Figure \ref{fig:CenterCrop}. This means that the center of the image will be cropped by one-fifth and two-fifths of its width and height. Then the image will be resized to the original size. 


\subsubsection{Flipping} \label{Flipping}
\begin{figure}
	 \centerline{\includegraphics[width=0.8\textwidth]{src/augmentations/flipping.png}}
	 {\caption{Examples of Flipping}\label{fig:flipping}}
\end{figure}
Flipping transformation is an image augmentation technique that flips an image horizontally or vertically. It involves reflecting the image along a vertical or horizontal axis, which results in a mirror image of the original image. 

The advantage of flipping is that it is a simple transformation that does not risk any loss of information. However, it is limited in terms of variety because it can only create mirrored versions of existing examples, which may not be sufficient for capturing the full range of variations in the dataset. Figure \ref{fig:flipping} illustrates the efffect of the transformations.

\subsubsection{Motion Blur} \label{MotionBlur}
Motion Blur transformation is an image augmentation technique that simulates Motion Blur in an image. The Motion Blur transformation works by applying a blur filter to the image, which simulates the effect of an object in motion. This is accomplished by convolving the image with a kernel function that creates the blur effect. Figure \ref{fig:motion_blur} illustrated the blur transformation. The size and direction of the kernel function can be varied to control the strength and direction of the Motion Blur effect. The blur limit in Motion Blur augmentation refers to the maximum amount of blur that can be applied to an image. This limit can be defined in terms of the size and direction of the kernel function used to apply the blur effect. 

\begin{figure}[!h]
	 \centerline{\includegraphics[width=0.8\textwidth]{src/augmentations/motion_blur.png}}
	 {\caption{Examples of Motion Blur with blur litmit of 30} \label{fig:motion_blur}}
\end{figure}


One advantage of Motion Blur transformation is that it can help to generate a more diverse set of training images for machine learning models. By varying the size and direction of the kernel function, it is possible to generate a large number of images with different types and strengths of Motion Blur. This can help improve the accuracy and generalization of the trained model.

While Motion Blur transformation can be useful for generating a diverse set of training images, it does have some significant drawbacks, particularly for object detection in aerial images of marine wildlife. A significant disadvantage of this technique is that it can reduce the image's sharpness and detail, making it more challenging to detect or recognize objects in the image. In the case of aerial images of marine wildlife, where shape and detail are essential for accurate detection, Motion Blur transformation can negatively impact the performance of the machine learning model. 

\subsubsection{Perspective} \label{Perspective}

Perspective transformation is an image augmentation technique that involves applying a geometric transformation to an image in order to alter its perspective. To perform perspective transformation, the algorithm first defines a set of four source points and a set of four destination points that define the corners of the original and transformed image, respectively. Then, it computes a transformation matrix that maps the source points to the destination points using techniques such as homography or projective geometry. An example of perspective transformation is the transformation of a squared image of a square into a trapezoid shape, then resize it to their original size. This transformation can create the illusion of a 3D perspective in a 2D image. 

Hyperparameter of Perspective includes random scale factor. The scale factor is used to determine the new points' distance from the original image corners during the transformation, which is calculated by sampling from a normal distribution with a mean of zero and a standard deviation set by the scale factor. A larger scale factor leads to a more pronounced transformation.  In the experiment, the scale factor will be tuned, the effect of Perspective transformation is shown in Figure \ref{fig:perspective}. 

\begin{figure}[!h]
	 \centerline{\includegraphics[width=0.8\textwidth]{src/augmentations/Perspective Transformation.png}}
	 {\caption{Example of Perspective with random scale between 0.1 and 0.2}\label{fig:perspective}}
\end{figure}

It should be noted that in the case of aerial images, perspective transformation may not be as effective. Aerial images are typically captured from a top-down perspective, where the camera is fixed in a position and oriented straight downward. As a result, the scene in the image is already in a consistent perspective, and applying a perspective transformation may not yield significant benefits.

\subsubsection{Rotation} \label{Rotation}

Rotation transformation is an image augmentation technique that involves rotating an image by a certain angle. The rotation can be performed in either a clockwise or counter-clockwise direction.


One advantage of rotation transformation is that it can increase the diversity of the training data without risking any information loss. This is particular useful for our dataset because marine wildlife aerial images often have objects in different orientations and angles. By applying rotation transformations, the model can learn to recognize the same object in various angles, which improves the model's robustness and accuracy in real-world scenarios. Additionally, rotation transformation is a simple and computationally efficient technique that can generate a wide variety of training examples. In this thesis, we apply a rotation transformation to our images within a range of -90 to 90 degrees. The impact of this transformation is demonstrated in Figure \ref{fig:rotation}.

\begin{figure}[!h]
	 \centerline{\includegraphics[width=0.8\textwidth]{src/augmentations/rotation.png}}
	 {\caption{Examples of Rotation}\label{fig:rotation}}
\end{figure}